\documentclass[a4paper,english]{scrartcl}
\usepackage[T1]{fontenc}
\usepackage[latin9]{inputenc}
\usepackage{amsbsy}

\makeatletter

\makeatother

\usepackage{babel}
\usepackage{listings}
\renewcommand{\lstlistingname}{Listing}

\begin{document}
\title{SEA3004F M4P1: surface elevation and the dispersion relation}
\date{Ocean and Atmosphere Dynamics}
\maketitle

\section{Introduction}
Download the scripts \texttt{wavelength.m} and \texttt{surface\_elevation.m} in a folder M4P1 on your PC.
This practical explores the surface elevation and the dispersion relation of oceanic gravity waves. The first exercise of each section will be explained in the tutorial and the others will be done on your own. You will submit a LiveScript with all the figures, comments, and answers.

\section{The surface elevation equation [20]}

Open the \texttt{surface\_elevation.m} file in the editor and look at the code. Read the comments at the start of the code for a description of what it does. Note that this function calls the wavelength.m function... i.e. a function within a function.

Execute the \texttt{surface\_elevation.m} function using an input amplitude of 0.5 m, a period of 16 s, and a water depth of 20 m:

\begin{lstlisting}[language=Matlab,basicstyle={\small},showstringspaces=false,frame=lines]
[eta,x,t]=surface_elevation(0.5,16,20);
\end{lstlisting}

type 'whos' to see which variables have been generated by the above command. Notice that eta is a matrix with 49 rows (corresponding to time) and 637 columns (corresponding to distance). We will plot this matrix and then extract some sections in time and space.

\subsection{[5]}

 Create a plot as outlined below (it is called a Hoevmoeller plot, which combines space and time). 

\begin{lstlisting}[language=Matlab,basicstyle={\small},showstringspaces=false,frame=lines]
figure
pcolor(x,t,eta);
shading interp;
title('Surface elevation');
ylabel('t (s)');
xlabel('x (m)');
c = colorbar;
ylabel(c,'eta (m)');
\end{lstlisting}

What do you think the slope of the lines of equal elevation represents?

\subsection{[5]}
 
Plot a time series of waves passing the x coordinate of 200 m. Label the figure appropriately.

\subsection{[10]}

Plot snapshots in time for times of 0 s and 4 s on the same figure, using the appropriate labels and legend. Explain in what direction are the orbital velocities of a water particle on the surface at x = 0 for these two time stamps.

\section{The dispersion relation [20]}

Open the \texttt{wavelength.m} file in the editor and look at the code. This function provides an iterative solution for the \textit{dispersion relation} from Lecture 4. The dispersion equation for gravity waves is the following:  

\[
\omega^2 = g k \tanh(kd)
\]

where:  
\begin{itemize}
    \item \( \omega \) is the wave angular frequency,
    \item \( g \) is the acceleration due to gravity,
    \item \( k \) is the wavenumber, and
    \item \( d \) is the water depth.
\end{itemize}

This is a transcendental equation, which means that it cannot be resolved by algebraic methods. The variable $k$ is both outside and inside a function that cannot be resolved. It requires a numerical approximation.
Go through the wavelength.m code and convince yourself that you understand what the code is doing. You will notice that the equation used in the code is slightly different from the one above. Here we make explicit the angular frequency and wavenumber using the period $T$ and the wavelength $L$:

\[
L_{out} = \frac{g T^2}{2\pi} \tanh \left(\frac{2\pi d}{L_{in}}\right)
\]

As L appears on both sides of the equation, we need an input wavelength \(L{in}\) which we initially guess as the deep-water value. We then check to see if \(L_{in}\) is close to the calculated \(L_{out}\) (within a tolerance). If it is close enough, then we stop, otherwise the code uses the newly calculated Lout as Lin for the next iteration. Once you are satisfied that you understand the code, test the wavelength function using an input wave period $T = 16$ s and depth $D = 50$ m.


\subsection{[10]}

Test the sensitivity of wavelength to water depth by computing the wavelength at depths ranging from 500 m to 10 m. One way to do this is using a for loop as follows:

\begin{lstlisting}[language=Matlab,basicstyle={\small},showstringspaces=false,frame=lines]
 T = 16;
 depths = 500:-10:10;
 L = zeros(size(depths));
 for d = 1:length(depths)
     L(d) = wavelength(T,depths(d));
 end
\end{lstlisting}

Generate a plot of depth vs. wavelength. Indicate your visual estimate of the maximum wavelength obtained for a wave of this period (T=16 s) and the water depth at which this is achieved.

\subsection{[10]}

Determine the depths at which a linear harmonic ocean surface gravity wave with a period of 16 s can be described as a \textit{shallow-water} and \textit{deep-water} wave, based on the guidelines of $d<L/20$ and $d>L/2$, respectively. (refer to the examples in Lecture 4).
These values will be used to answer the following questions. Write the code to: 
 \begin{enumerate}
     \item  Calculate the \textit{deep-water} wavelength using the deep-water dispersion relation and compare this value with the wavelength calculated by the wavelength.m function. Use the minimum deep-water depth limit calculated above. Calculate the percentage error we make by assuming the simplified deep-water dispersion relation at this depth. [5]
     \item Calculate the \textit{shallow-water} wavelength using the simplified dispersion relation. Use the maximum shallow-water depth limit calculated above. Show the percentage error we make by assuming the simplified dispersion relation at this depth. [5] 
 \end{enumerate}
\end{document}